\setchapterpreamble[u]{\margintoc}
\chapter{Euler Tour Trees}
\labch{Introduzione}

In questa tesi esploreremo gli Euler Tour Trees 
(ETT), una struttura dati dinamica introdotta da 
Henzinger e King \cite{HenzingerKing1999}, 
progettata per mantenere informazioni su foreste 
soggette a modifiche strutturali.

L'obiettivo è supportare operazioni di: 
\begin{itemize}
    \item \textbf{link(u,v)}: aggiungere un arco $(u,v)$ unendo due alberi.
    \item \textbf{cut(u,v)}: rimuovere l'arco $(u,v)$ dividendo l'albero in due.
    \item \textbf{isConnected(u,v)}: verificare se $u$ e $v$ appartengono allo stesso albero.
\end{itemize}

A differenza di altre strutture che operano direttamente sui nodi,
gli ETT si basano sulla linearizzazione dell'albero 
tramite cicli di Eulero. Questa rappresentazione
permette di ridurre complessi problemi di connettività 
a semplici manipolazioni di sequenze.

\begin{definition}
Un ciclo euleriano è un percorso in un grafo che inizia e finisce nello stesso nodo passando esattamente una volta per ogni arco.
\end{definition}

\begin{kaobox}[frametitle=Note] 
Poiché gli alberi non hanno veri cicli euleriani, consideriamo ogni arco $e=(u,v)$ come una coppia di archi diretti $(u,v)$ e $(v,u)$.
\end{kaobox}

L'idea della stuttura dati è quella di creare il ciclo euleriano visitando l'albero tramite una ricerca in profondità (DFS) tenendo traccia degli archi attraverso i quali si prosegue nella visita

\begin{marginfigure}
    \centering
    \includegraphics[width=\marginparwidth]{makeTour.png}
    \caption[Ciclo Euleriano]{Il Ciclo euleriano che si forma è (A,B), (B,D), (D,B), (B,E), (E,B), (B,A), (A,C), (C,A)}
    \label{fig:Tour}
\end{marginfigure}

\begin{algorithm}[H]
    \caption{Creazione Tour Euleriano }
    \label{alg:make_ETT}
    \SetAlgoLined
    \DontPrintSemicolon
    
    \KwIn{Nodo corrente $u$, Lista $L$}
    \KwOut{Cammino euleriano in $L$}
    
    \BlankLine 

    \ForEach{figlio $v$ di $u$}{
        $L.append( (u, v) )$ \;
        makeTour($v, L$)\;
        $L.append( (v, u) )$ \;
    }

\end{algorithm}

Il ciclo euleriano è strettamente connesso con la struttra dell'albero, ad esempio si osserva che la sequanza (A,B), \textbf{(B,D), (D,B), (B,E), (E,B)}, (B,A), (A,C), (C,A) rappresenta il sottoalbero radicato in B

\begin{kaobox}[frametitle = Claim]
    Sia $S$ la sequenza di archi che descrive 
    il tour euleriamo per ogni nodo $v$ il 
    sottoalbero radicato in $v$ è rappresenta dall'intervallo 
    contiguo va dal primo arco che ha $(v,\_)$ fino all'ultimo arco che contiene $(\_,v)$    
\end{kaobox}\setchapterpreamble[u]{\margintoc}

\subsection{Operazioni}

\subsubsection{isConnected}
L'operazione \textbf{isConnected(u,v)} è la più semplice poichè è sufficente verificare che i due nodi $u$ e $v$ si trovino nello stesso ciclo euleriano, per farlo basta verificare che il primo elemento della sequanza sia il medesimo per i due nodi.

\subsubsection{Rerooting}
Prima di introdurre le altre operazioni introduciamo,
un'operazione ausiliaria \textbf{rerooting(T,v)}, 
tale operazione ha il compito di rendere un nodo 
interno del ciclo euleriano il primo nodo del ciclo
, per farlo seguiamo i seguenti passaggi:

\begin{enumerate}
    \item si cerca una occorrenza che ha $(v,c)$ (dove $c$ è un figlio casuale di v) nel ciclo euleriano, e si separa la sequanza in A che conterrà tutto ciò che c'era prima di tale arco e B che conterrà l'arco $(v,c)$ e tutto ciò che c'è dopo
    \item riconcateno $B + A + (c,v)$
\end{enumerate}

\begin{algorithm}[H]
    \caption{Rerooting}
    \label{alg:rerooting}
    \SetAlgoLined
    \DontPrintSemicolon
    
    \KwIn{ETT $L$, nodo $v$}
    \KwOut{Il nuovo ETT con radice $v$}
    
    \BlankLine 

    \If{$L[0].source == v$}{
        \Return $L$
    }
    
    $A \leftarrow \emptyset; B \leftarrow \emptyset; isFound \leftarrow False; c \leftarrow NULL$\;
    

    \For{$i \leftarrow 0$ \KwTo $len(L) - 1$}{
        \tcp{Cerchiamo il primo arco che parte da v}
        \If{$L[i].source == v$ \textbf{and} $\neg isFound$}{
            $isFound \leftarrow True$\;
            $c \leftarrow L[i].target$\;
        }

        \If{$\neg isFound$}{
            $A \leftarrow A + L[i]$\;
        }
        \Else{
            $B \leftarrow B + L[i]$\;
        }
    }

    \Return $B + A + (c,v)$
\end{algorithm}

\subsubsection{Link}
L'operazione \textbf{link(u,v)} permette di unire due alberi distinti aggiungendo un arco tra il nodo u (appartenente al primo albero) e il nodo v (appartenente al secondo), per farlo seguiamo i seguenti passaggi:
\begin{enumerate}
    \item A = rerooting(u)
    \item B = rerooting(v) 
    \item concateno  $A + (u,v) + B + (v,u)$ 
\end{enumerate}

\begin{algorithm}[H]
    \caption{Link}
    \label{alg:link}
    \SetAlgoLined
    \DontPrintSemicolon
    
    \KwIn{ETT $L$, ETT $C$, nodo $u$, nodo $v$}
    \KwOut{l'ETT del nuovo albero}
    
    \BlankLine 

    $A \leftarrow rerooting(L,u)$\;
    $B \leftarrow rerooting(C,v)$\;

    \Return $A + (u,v) + B + (v,u)$
    
\end{algorithm}

\subsubsection{Cut}
L'operazione \textbf{cut(u,v)} permette di separare un albero in due alberi distinti, eliminando l'arco $(u,v)$ crando due nuovi alberi radicati rispettivamente in $u$ e $v$, per farlo seguiamo i seguenti passaggi:
\begin{enumerate}
    \item rerooting(T,v)
    \item si cerca la sequanza $(u,v)$ e $(v,u)$ e si separea in A che conterrà tutto ciò che c'era prima di $(u,v)$, B che conterrà $(u,v)$ e $(v,u)$ e tutto ciò che c'è tra i due e C che contiene ciò che c'è dopo $(v,u)$
    \item si eliminano il primo e l'ultimo elemento in $B$
    \item si restituisce $A + C$ e $B$
\end{enumerate}

\begin{algorithm}[H]
    \caption{Cut}
    \label{alg:cut}
    \SetAlgoLined
    \DontPrintSemicolon
    
    \KwIn{ETT $L$, nodo $u$, nodo $v$}
    \KwOut{Due ETT: $L_{rimanente}$ e $L_{staccato}$}
    
    \BlankLine 

    $L \leftarrow \text{rerooting}(L, v)$\;
    
    $A \leftarrow \emptyset; B \leftarrow \emptyset; C \leftarrow \emptyset$\;
    

    \tcp{Trovo il primo arco in cui compare (u,\_)}

    \For{$i \leftarrow 0$ \KwTo $\text{len}(L)-1$}{

        \If{$L[i].source == u$}{
            \textbf{break}\;
        }

        $A \leftarrow A + L[i]$\;

    }

    \tcp{Trova l'ultima arco in cui compare (\_,u)}

    \For{$j \leftarrow \text{len}(L)-1$ \KwTo $0$}{
        \If{$L[j].target == u$}{
            \textbf{break}\;
        }

        $C \leftarrow L[j] + C$\;
    }


    \tcp{Costruisce la sequenza B tra le due occorrenze}

    \For{$k \leftarrow i$ \KwTo $j$}{
        $B \leftarrow B + L[k]$\;
    }

    \text{elimino il primo e l'ultimo elemento di } $B$\;

    \Return $(A + C)$, $B$ 
\end{algorithm}

 \begin{remark}
    Tutte queste implementazioni delle operazioni però hanno costo proporzionale alla dimensione del ciclo euleriano dunque il loro costo sara $O(n)$
 \end{remark}

L'obbiettivo quindi sara quello di introdurre una struttra dati che ci permetterà di eseguire queste operazioni in maniera molto piu efficente
